\documentclass[11pt]{article}
\usepackage[margin=0.95in]{geometry}
\usepackage[T1]{fontenc}
\usepackage{lmodern}
\usepackage[hidelinks]{hyperref}
\usepackage{booktabs}
\usepackage{microtype}
\newcommand{\code}[1]{\texttt{\detokenize{#1}}}

\title{Fuzz Campaign Physical-Core Parallelism}
\author{Tactix Development Project}
\date{26 August 2026}

\begin{document}
\emergencystretch=2em
\maketitle

\section*{Outcome}
The active consolidation fuzz campaign is correctness-clean so far but was
submitted at the wrong parallel width.  Claude Code's ownership audit reports
six generated/stress legs as independent \code{-c 2} jobs with two workers
each.  On the cluster's SMT-2 nodes this is approximately one physical core per
leg, or about six physical cores at peak.  The campaign has no arrays or hidden
throttle and never achieved the intended 50-physical-core parallelism.

This affects wall time, not seeds, case count, coverage, or oracle semantics.
At discovery the two remaining legs were already at 8,700/10,000 and
9,000/10,000 cases with zero failures.  They are retained to completion rather
than cancelled and repeated.  No new fuzz campaign was submitted.

The generic fuzz wrapper is corrected so its default ``50'' means 50 physical
cores and 50 workers, rather than 50 logical CPUs and 50 workers sharing 25
physical cores.  The Development Plan now records the same contract for future
Claude-owned generated fuzz legs.

\section{Observed campaign shape}
Claude reports prep jobs \code{1648187}/\code{1648211}, generated/stress legs
\code{1648214}--\code{1648218}, and reach job \code{1648219}.  Each generated
leg requests two logical CPUs and 16 GiB, and launches two workers.  Four legs
are terminal clean; Congruence-stress and Sweep-stress remain the two original
in-flight legs.  Claude, not Codex, inspected these job records and fuzz
artifacts, consistent with Development Plan \S8.2.

The prior \code{tools/fuzz/slurm.sh} default requested
\code{--cpus-per-task=50} with \code{--jobs 50}.  Because the nodes expose two
hardware threads per physical core, that generic shape reserved only 25
physical cores.  The active custom campaign did not use even that wrapper and
instead used its two-worker leg shape.

\section{Validated future contract}
Allocation-only job \code{1648629} ran no solver and no fuzz case.  It completed
0:0 and reports:

\begin{center}
\begin{tabular}{ll}
\toprule
Resource & Observed value \\
\midrule
Logical CPUs & 100 \\
Scheduler memory & 400,000 MiB \\
Allowed CPU mask & \code{0-49,96-145} \\
Physical cores represented & 50 \\
\bottomrule
\end{tabular}
\end{center}

The mask contains cores 0--49 and their SMT siblings 96--145, proving that 100
logical CPUs reserve 50 physical cores without relying on a scheduler hint.
Memory is 50 times the partition's effective 8,000 MiB per physical core.
Future generated fuzz legs launch 50 workers, leaving one worker per physical
core while the allocation reserves both siblings.

The wrapper now accepts \code{--physical-cores P}, derives
\code{cpus-per-task=2P}, \code{mem=8000P MiB}, and \code{jobs=P}, and rejects
values outside 1--96.  A mocked submission test verifies the default
100/400,000/50 shape; a three-core control verifies 6/24,000/3 scaling; zero is
rejected before submission.  Shell syntax validation passes.

Preparation, small fixed-corpus work, or intrinsically narrower reachability
probes may use less only with a recorded reason.  A generated fuzz leg may not
silently fall back to one or two workers.  This is parallelism inside one fuzz
campaign, not the three-wave repetition used for decision-grade performance
timing.

\section{Next action}
Let Claude finish and audit the two nearly complete legs of the existing single
campaign.  A terminal PASS authorizes the already-complete consolidation gate
to proceed to its own final report, commit, and push; a terminal failure enters
the mandatory bounded Claude repair workflow.  The corrected 50-physical-core
wrapper applies to the next required fuzz campaign and does not retroactively
invalidate the current campaign's completed coverage.

\end{document}
